%=====================================================================
\chapter{Python Quick Reference}\label{app:python}
%=====================================================================
\normalfigures

Everything here appears somewhere in the handout. This appendix collects it so you
can work without paging back.

\section{Setting Up}

\begin{lstlisting}[language=bash]
python -m venv .venv                      # one environment per project
.\.venv\Scripts\Activate.ps1              # Windows PowerShell
pip install pandas numpy scikit-learn matplotlib seaborn scipy statsmodels
pip freeze > requirements.txt             # pin versions -- Chapter 29
\end{lstlisting}

\section{pandas --- Loading and Inspecting}

\begin{lstlisting}[language=Python]
import pandas as pd, numpy as np

df = pd.read_csv("data.csv", parse_dates=["timestamp"])
df = pd.read_parquet("data.parquet")      # faster, keeps dtypes

df.shape            # (rows, columns)
df.dtypes           # types -- check categories are not stored as numbers
df.head(10)
df.describe()       # numeric summary: count, mean, sd, quartiles
df.describe(include="object")             # categorical summary
df.info(memory_usage="deep")

df.isna().mean().sort_values(ascending=False)   # % missing per column
df.duplicated().sum()
df.nunique().sort_values()                # few distinct + numeric => a category
\end{lstlisting}

\section{Selecting, Filtering, Grouping}

\begin{lstlisting}[language=Python]
df["col"]                    # one column (Series)
df[["a", "b"]]               # several columns (DataFrame)
df.loc[df.age > 30, ["a"]]   # filter rows by condition, select columns
df.iloc[0:5, 0:3]            # by position

df.query("age > 30 and programme == 'BSIT'")

df.groupby("programme").agg(
    n=("student_id", "size"),
    mean_mark=("mark", "mean"),
    pass_rate=("passed", "mean"),
)

df.pivot_table(index="programme", columns="year", values="mark", aggfunc="mean")
pd.crosstab(df.programme, df.grade, normalize="index")   # row proportions
\end{lstlisting}

\section{Cleaning (Chapter 7)}

\begin{lstlisting}[language=Python]
df = df.drop_duplicates(subset=["student_id", "term"])    # deduplicate FIRST

df["prior"] = df["prior"].fillna(df["prior"].median())
df["prior_was_missing"] = df["prior"].isna().astype(int)  # MNAR: encode, don't hide

df["temp"] = df["temp"].ffill(limit=6)                    # cap forward fill

df = df.replace({"age": {0: np.nan, -1: np.nan}})         # placeholder values

# joins -- ALWAYS compare row counts before and after
before = len(df)
df = df.merge(vle, on="student_id", how="left")           # left, not inner
assert len(df) == before, "join changed row count -- key is not unique"

# outliers: flag, do not silently delete
q1, q3 = df.dur.quantile([.25, .75])
iqr = q3 - q1
df["outlier"] = ~df.dur.between(q1 - 1.5*iqr, q3 + 1.5*iqr)
\end{lstlisting}

\section{scikit-learn --- The Standard Workflow}

\begin{lstlisting}[language=Python]
from sklearn.model_selection import train_test_split, StratifiedKFold, cross_val_score
from sklearn.pipeline import Pipeline, make_pipeline
from sklearn.compose import ColumnTransformer
from sklearn.impute import SimpleImputer
from sklearn.preprocessing import StandardScaler, OneHotEncoder

# 1. SPLIT FIRST -- before anything is learned (Chapter 7)
X_tr, X_te, y_tr, y_te = train_test_split(
    X, y, test_size=0.2, stratify=y, random_state=42)

# 2. Pipeline: makes leakage structurally impossible
prep = ColumnTransformer([
    ("num", make_pipeline(SimpleImputer(strategy="median", add_indicator=True),
                          StandardScaler()), numeric_cols),
    ("cat", OneHotEncoder(handle_unknown="ignore"), categorical_cols),
])
pipe = make_pipeline(prep, model)

pipe.fit(X_tr, y_tr)          # fit on train
pipe.predict(X_te)            # transform-only on test -- never fit_transform
\end{lstlisting}

\subsection*{Models by task}

\begin{lstlisting}[language=Python]
# --- regression (Chapter 12) ---
from sklearn.linear_model import LinearRegression, RidgeCV, LassoCV
from sklearn.ensemble import RandomForestRegressor, HistGradientBoostingRegressor

# --- classification (Chapters 13, 15) ---
from sklearn.linear_model import LogisticRegression
from sklearn.neighbors import KNeighborsClassifier
from sklearn.naive_bayes import GaussianNB, MultinomialNB
from sklearn.tree import DecisionTreeClassifier, export_text
from sklearn.svm import SVC
from sklearn.ensemble import RandomForestClassifier, HistGradientBoostingClassifier

# --- unsupervised (Chapter 16) ---
from sklearn.cluster import KMeans, DBSCAN, AgglomerativeClustering
from sklearn.decomposition import PCA
from sklearn.ensemble import IsolationForest

# --- baselines. ALWAYS. (Chapter 11) ---
from sklearn.dummy import DummyClassifier, DummyRegressor
\end{lstlisting}

\subsection*{Evaluation (Chapter 14)}

\begin{lstlisting}[language=Python]
from sklearn.metrics import (confusion_matrix, classification_report,
                             roc_auc_score, average_precision_score,
                             precision_recall_curve, mean_absolute_error,
                             mean_squared_error, r2_score)

tn, fp, fn, tp = confusion_matrix(y_te, y_pred).ravel()
print(classification_report(y_te, y_pred, digits=3))

proba = pipe.predict_proba(X_te)[:, 1]
average_precision_score(y_te, proba)      # PR-AUC: use when positives are rare
roc_auc_score(y_te, proba)                # ROC-AUC: flatters rare-event detectors

# choose the threshold from costs, not from 0.5
prec, rec, thr = precision_recall_curve(y_te, proba)

# cross-validation -- match the scheme to the data
StratifiedKFold(5, shuffle=True, random_state=0)   # imbalanced classes
from sklearn.model_selection import GroupKFold, TimeSeriesSplit
GroupKFold(5)          # many rows per student/device
TimeSeriesSplit(5)     # anything temporal -- never random
\end{lstlisting}

\section{Time Series (Chapter 19)}

\begin{lstlisting}[language=Python]
s = df.set_index("date")["value"].asfreq("D")

s.shift(1)                       # PAST. shift(-1) is the future -- a leak.
s.shift(1).rolling(7).mean()     # shift FIRST, then roll: trailing window only
s.rolling(7, center=True)        # <-- LEAK. never use center=True for features.

s.diff(1); s.pct_change(7)
s.resample("W").sum()

from statsmodels.tsa.seasonal import seasonal_decompose
seasonal_decompose(s, model="additive", period=7).plot()
\end{lstlisting}

\section{Plotting (Chapter 8)}

\begin{lstlisting}[language=Python]
import matplotlib.pyplot as plt, seaborn as sns

df.hist(bins=30, figsize=(11, 7))                       # univariate
sns.boxplot(data=df, x="programme", y="mark")           # numeric by category
sns.scatterplot(data=df, x="attendance", y="mark", hue="programme")
sns.heatmap(df.corr(numeric_only=True), annot=True, fmt=".2f",
            cmap="RdBu_r", center=0)                    # diverging for signed
sns.lmplot(data=df, x="hours", y="mark", hue="programme")   # per group -- Simpson's

plt.tight_layout(); plt.savefig("fig.png", dpi=150, bbox_inches="tight")
\end{lstlisting}

\section{Reproducibility (Chapter 29)}

\begin{lstlisting}[language=Python]
import random, os
SEED = 42
random.seed(SEED); np.random.seed(SEED)
os.environ["PYTHONHASHSEED"] = str(SEED)
# every sklearn estimator: random_state=SEED
\end{lstlisting}

\begin{alertbox}[The Five Lines That Cause Most Bugs]
\begin{tightnum}
  \item \texttt{scaler.fit(X)} before \texttt{train\_test\_split} --- leakage.
  \item \texttt{.fit\_transform(X\_test)} --- leakage.
  \item \texttt{rolling(n, center=True)} in a feature --- leakage from the future.
  \item \texttt{train\_test\_split} on time-series data --- trains on the future.
  \item \texttt{accuracy\_score} on imbalanced data --- meaningless.
\end{tightnum}
\end{alertbox}
